\documentclass[../DoAn.tex]{subfiles}
\begin{document}

Năm chương trước đã trình bày toàn bộ quá trình thực hiện đồ án, từ đặt vấn
đề, khảo sát hiện trạng, lựa chọn công nghệ, thiết kế và hiện thực hệ thống
\break ALIBARBIE, đến các đóng góp kỹ thuật đặc thù. Chương này tổng kết kết quả đạt
được so với mục tiêu đề ra ở mục \ref{section:1.2}, đối chiếu hệ thống với hai
nền tảng thương mại điện tử phổ biến là Haravan và Shopify, nhìn nhận các hạn
chế còn tồn tại, và đề xuất hướng phát triển trong tương lai.

\section{Kết luận}
\label{section:6.1}

\subsection*{So sánh với các trang web thương mại điện tử hiện có đại diện tiêu biểu là  Haravan và Shopify}

Haravan đại diện cho nhóm nền tảng SaaS thương mại điện tử trong nước, còn
Shopify đại diện cho nhóm nền tảng quốc tế có khả năng tùy biến sâu hơn qua hệ
thống ứng dụng mở rộng (bảng \ref{tab:comparison}). Đối chiếu với hai nền tảng
này theo ba tiêu chí đã được xác định là hạn chế chung của thị trường ở mục
\ref{section:2.1}, ALIBARBIE đạt được ba kết quả vượt trội cụ thể. Thứ nhất,
quy trình hoàn trả hàng của ALIBARBIE được mô hình hóa thành máy trạng thái
hữu hạn với bảy trạng thái (mục \ref{subsection:5.1.2}), trong khi Haravan và
Shopify ở cấu hình mặc định chỉ hỗ trợ hai trạng thái chấp nhận hoặc từ chối,
muốn mở rộng vòng đời chi tiết hơn phải thuê thêm ứng dụng của bên thứ ba qua
App Store. Thứ hai, ALIBARBIE tích hợp sẵn nhật ký tác động hệ thống với tham chiếu đa
hình (mục \ref{subsection:5.2.2}), còn hai nền tảng đối chứng không cung cấp
tính năng này trong gói dịch vụ cơ bản. Thứ ba, ALIBARBIE cho phép cấu hình
phí vận chuyển đến cấp xã/phường (mục \ref{subsection:5.3.2}), trong khi
Haravan và Shopify chỉ cấu hình theo vùng hoặc theo biểu giá của đối tác vận
chuyển tích hợp, không tùy biến đến cấp địa bàn nhỏ nhất.

Tuy vậy, xét trên tổng thể, Haravan và Shopify vẫn vượt trội ALIBARBIE ở các
mặt mà một đồ án tốt nghiệp không thể đạt được trong phạm vi và thời gian thực
hiện. Hai nền tảng này có hệ sinh thái ứng dụng mở rộng phong phú, đã được
kiểm chứng khả năng chịu tải ở quy mô hàng triệu cửa hàng thực tế, hỗ trợ đa
cổng thanh toán và đa ngôn ngữ, có sẵn ứng dụng di động cho cả người mua và
người quản lý, và được vận hành bởi đội ngũ kỹ thuật chuyên biệt nên doanh
nghiệp sử dụng không phải tự duy trì hạ tầng. ALIBARBIE bù lại các hạn chế này
bằng quyền sở hữu hoàn toàn dữ liệu và mã nguồn, không phát sinh chi phí thuê
bao định kỳ, và khả năng tùy chỉnh nghiệp vụ không bị giới hạn bởi nền tảng
của bên thứ ba.

\subsection*{Kết quả đạt được so với mục tiêu}

Đối chiếu với các mục tiêu đề ra ở mục \ref{section:1.2}, đồ án đã hiện thực
đầy đủ các chức năng dự kiến: quản lý sản phẩm theo biến thể và danh mục;
giỏ hàng, đặt hàng và thanh toán trực tuyến qua VNPay kết hợp thanh toán khi
nhận hàng; quản lý vòng đời đơn hàng và quy trình hoàn trả hàng đa trạng
thái; cấu hình phí vận chuyển theo địa bàn; quản lý mã giảm giá, danh sách
yêu thích và blog; phân quyền nhân viên theo vai trò kết hợp nhật ký kiểm
toán; và hỗ trợ khách hàng qua chatbot thời gian thực. Toàn bộ 21 trường hợp
kiểm thử trọng điểm đạt kết quả đúng (mục \ref{section:4.4}), và kết quả
triển khai thử nghiệm xác nhận hệ thống đáp ứng các yêu cầu phi chức năng đã
đặt ra ở mục \ref{section:2.4} về thời gian phản hồi và tính tương thích
trình duyệt.

Bên cạnh các chức năng đã hoàn thành, đồ án còn tồn tại một số hạn chế. Hệ
thống mới được triển khai và kiểm thử trên môi trường cục bộ (mục
\ref{section:4.5}), chưa được vận hành trên hạ tầng cloud thực tế nên chưa
có cơ sở đánh giá khả năng chịu tải khi nhiều người dùng truy cập đồng thời.
Phạm vi kiểm thử tập trung vào 21 trường hợp trọng điểm và các trường hợp bổ
sung trong phụ lục, chưa bao gồm kiểm thử tải và kiểm thử an ninh thâm nhập
toàn diện. Hệ thống chỉ cung cấp giao diện web đáp ứng, chưa có ứng dụng di
động riêng cho khách hàng và nhân viên. Hình thức thanh toán trực tuyến hiện
chỉ giới hạn ở cổng VNPay, chưa hỗ trợ các ví điện tử phổ biến khác. Cuối
cùng, hệ thống gợi ý sản phẩm cho khách hàng vẫn dựa trên tìm kiếm và lọc thủ
công theo danh mục, chưa có cơ chế gợi ý cá nhân hóa dựa trên hành vi mua
hàng.

\subsection*{Đóng góp nổi bật và bài học kinh nghiệm}

Ba đóng góp kỹ thuật đặc thù của đồ án, đã được trình bày chi tiết ở Chương 5,
là: mô hình máy trạng thái hữu hạn cho quy trình hoàn trả hàng với kiểm soát
chuyển đổi và phân quyền chặt chẽ theo từng bước; mô hình RBAC hai tầng kết
hợp nhật ký kiểm toán tham chiếu đa hình, cho phép truy vết mọi hành động của
nhân viên; và cơ chế cấu hình phí vận chuyển phân cấp đến cấp xã/phường dựa
trên cấu trúc tài liệu lồng nhau của MongoDB. Bên cạnh đó, cơ chế ký số
HMAC-SHA512 cho luồng thanh toán VNPay (mục \ref{subsection:5.4.2}) đảm bảo
tính toàn vẹn của giao dịch tài chính, là một đóng góp bổ trợ quan trọng cho
độ an toàn của hệ thống.

Quá trình thực hiện đồ án mang lại một số bài học kinh nghiệm. Việc thiết kế
mô hình trạng thái rõ ràng ngay từ giai đoạn phân tích yêu cầu giúp tránh
được việc chỉnh sửa lớn về sau khi nghiệp vụ hoàn trả hàng phát sinh thêm
bước xử lý. Kiểm soát phân quyền chỉ đáng tin cậy khi được thực hiện ở tầng
server; việc ẩn hiện thành phần giao diện theo vai trò chỉ có giá trị về
trải nghiệm người dùng và không thể thay thế cho việc kiểm tra quyền tại
từng API. Mô hình tài liệu lồng nhau của MongoDB phù hợp hơn so với chuẩn
hóa quan hệ khi dữ liệu có cấu trúc phân cấp cố định và được truy vấn theo
chuỗi cha-con, như trường hợp cấu hình phí vận chuyển theo tỉnh/thành, quận/
huyện và xã/phường. Cuối cùng, khi tích hợp cổng thanh toán của bên thứ ba,
việc tuân thủ chặt chẽ tài liệu kỹ thuật chính thức và chủ động xác thực chữ
ký ở cả hai chiều gửi và nhận là điều kiện bắt buộc để tránh các lỗ hổng bảo
mật nghiêm trọng.

\section{Hướng phát triển}
\label{section:6.2}

Trước khi mở rộng thêm chức năng mới, đồ án xác định một số công việc cần
thiết để hoàn thiện các chức năng đã xây dựng. Hệ thống cần được triển khai
lên hạ tầng cloud thay cho môi trường cục bộ hiện tại, kèm theo quy trình
tích hợp và triển khai liên tục (CI/CD) để giảm thời gian đưa các thay đổi
vào vận hành. Phạm vi kiểm thử cần được mở rộng sang kiểm thử tải để xác định
ngưỡng chịu tải đồng thời, và kiểm thử an ninh thâm nhập để phát hiện sớm các
lỗ hổng trước khi vận hành thực tế. Bộ kiểm thử tự động (unit test và
integration test) cũng cần được bổ sung cho các module hiện chỉ được kiểm
thử thủ công, đồng thời hệ thống cần có cơ chế giám sát lỗi tập trung để phát
hiện và cảnh báo sự cố kịp thời.

Sau khi hoàn thiện các chức năng hiện có, đồ án đề xuất bốn hướng phát triển
mới. Hướng thứ nhất là xây dựng ứng dụng di động bằng React Native, tái sử
dụng phần lớn các API đã hiện thực ở phía server, cho phép khách hàng mua
sắm và nhân viên xử lý đơn hàng, yêu cầu hoàn trả mọi lúc mọi nơi thay vì chỉ
qua giao diện web đáp ứng. Hướng thứ hai là tích hợp công nghệ trí tuệ nhân
tạo để gợi ý sản phẩm cá nhân hóa dựa trên lịch sử mua hàng và hành vi duyệt
web của khách hàng, có thể áp dụng theo hướng lọc dựa trên nội dung sản phẩm
hoặc lọc cộng tác giữa các khách hàng có hành vi tương tự, nhằm tăng tỷ lệ
chuyển đổi so với cơ chế tìm kiếm và lọc thủ công hiện tại. Hướng thứ ba là
mở rộng cổng thanh toán trực tuyến, bổ sung ví điện tử MoMo bên cạnh VNPay
hiện có, áp dụng lại cơ chế ký số HMAC đã được kiểm chứng ở mục
\ref{subsection:5.4.2} để đảm bảo an toàn giao dịch trên cổng thanh toán mới.
Hướng thứ tư là xây dựng phân hệ thống kê và báo cáo cho quản trị viên với
biểu đồ doanh thu và sản phẩm bán chạy theo thời gian thực, tận dụng dữ liệu
đơn hàng và nhật ký kiểm toán đã được thu thập sẵn để hỗ trợ ra quyết định
kinh doanh.

\end{document}
